\documentclass[12pt,a4paper]{article}
        \makeatletter
        \def\input@path{{../}}
        \makeatother
        \usepackage{almeria-fichas}

        \FichaMeta{Sesion 17}{Seleccionar la informacion de la guia}{Bloque 4 · Producto final}{Filtrar mapas, calculos y graficas para elegir solo la informacion mas clara, precisa y relevante para la guia final.}{Procedimientos}{Analizar, Evaluar, Crear}

        \begin{document}

        \FichaCabecera

        \begin{materialesbox}
        \begin{itemize}
    \item Ficha impresa
    \item Lapiz
    \item Borradores del proyecto
    \item Rubrica o lista de comprobacion del equipo
\end{itemize}
        \end{materialesbox}

        \begin{pasosbox}
        \begin{enumerate}
    \item Reune todos los materiales del proyecto.
    \item Elige primero lo que este bien calculado y bien explicado.
    \item Descarta lo repetido o poco claro.
    \item Anota por escrito cada decision del equipo.
\end{enumerate}
        \end{pasosbox}

        \begin{recordatoriobox}
        \begin{itemize}
    \item No todo lo hecho tiene que entrar en la guia final.
    \item Una seleccion buena combina claridad, correccion y utilidad.
    \item Hay que justificar por que se incluye o se descarta algo.
\end{itemize}
        \end{recordatoriobox}

        \begin{apoyobox}
        \begin{itemize}
    \item Puedes usar una tabla con dos columnas: entra / no entra.
    \item Marca con una estrella lo imprescindible.
    \item Si algo no tiene fuente, unidad o explicacion, revisalo antes de incluirlo.
\end{itemize}
        \end{apoyobox}

        \BloomSection{analizar}

\BloomExercise{analizar}{1}{Candidatos a la guia}{Analiza una lista de ocho posibles datos y decide cuales son fiables, claros y utiles para una guia matematica.}{8}

\BloomExercise{analizar}{2}{Informacion repetida}{Analiza dos materiales que dicen casi lo mismo. Explica cual conservarias y cual resumirias.}{7}

\BloomSection{evaluar}

\BloomExercise{evaluar}{1}{Razonar una seleccion}{Evalua si una grafica sin unidades deberia entrar en la guia. Justifica con criterios matematicos.}{7}

\BloomExercise{evaluar}{2}{Orden de secciones}{Evalua que estructura seria mejor: empezar por mapa, por datos o por rutas. Explica tu decision.}{7}

\BloomSection{crear}

\BloomExercise{crear}{1}{Indice provisional}{Crea un indice provisional de la guia con entre 4 y 6 apartados.}{8}

\BloomExercise{crear}{2}{Portada matematica}{Diseña la idea de una portada que incluya titulo, subtitulo, una estadistica y un mini mapa con tres puntos.}{8}

        \begin{evidenciabox}
        \begin{itemize}
    \item Aplicar criterios para elegir contenidos.
    \item Justificar decisiones de inclusion y descarte.
    \item Proponer una estructura inicial coherente.
\end{itemize}
        \end{evidenciabox}

        \end{document}
